

\section{Database Management Systems}
%% ═════════════════════════════════════════════════════════════════════════════

\begin{frame}{\textcolor{brown}{What is a DBMS?}}
  \begin{block}{Definition}
    A \textbf{Database Management System (DBMS)} is software that enables users to
    \textbf{create, define, maintain, and control} access to a database.
    It acts as the interface between the user and the stored data, ensuring data
    remains consistently organised and accessible.
  \end{block}

  \vspace{0.7em}
  \begin{columns}[T]
    \column{0.48\textwidth}
      \begin{block}{\faTable\; Relational (SQL)}
        \begin{itemize}
          \item Data stored in structured \textbf{tables} (rows \& columns)
          \item Tables linked via \textbf{keys}
          \item[\faAngleRight] PostgreSQL, MySQL, Oracle
        \end{itemize}
      \end{block}

    \column{0.48\textwidth}
      \begin{block}{\faProjectDiagram\; Non-Relational (NoSQL)}
        \begin{itemize}
          \item Flexible formats: documents, key-value pairs, graphs
          \item Suited for unstructured or widely varied data
          \item[\faAngleRight] MongoDB, Redis, Firebase, Supabase
        \end{itemize}
      \end{block}
  \end{columns}

  \vspace{0.5em}
  \centering\small
  This seminar focuses on \textbf{Relational Databases} and the language used to interact with them — \textbf{SQL}.
\end{frame}
